\documentclass[11pt]{article}

\usepackage[T1]{fontenc}
\usepackage{lmodern}
\usepackage[margin=1in]{geometry}
\usepackage{amsmath,amssymb,amsthm,mathtools}
\usepackage{bm}
\usepackage{booktabs}
\usepackage{enumitem}
\usepackage{microtype}
\usepackage{xcolor}
\usepackage[numbers,sort&compress]{natbib}
\usepackage[colorlinks=true,linkcolor=blue!55!black,citecolor=blue!55!black,urlcolor=blue!55!black]{hyperref}

\newtheorem{theorem}{Theorem}[section]
\newtheorem{lemma}[theorem]{Lemma}
\newtheorem{proposition}[theorem]{Proposition}
\newtheorem{corollary}[theorem]{Corollary}
\newtheorem{definition}[theorem]{Definition}
\theoremstyle{remark}
\newtheorem{remark}[theorem]{Remark}

\newcommand{\R}{\mathbb R}
\newcommand{\E}{\mathbb E}
\newcommand{\Pp}{\mathbb P}
\newcommand{\1}{\bm 1}
\newcommand{\eps}{\varepsilon}
\newcommand{\cP}{\mathsf C_{\mathrm P}}
\newcommand{\Rq}{\mathcal R}
\newcommand{\TV}{\mathsf{TV}}
\newcommand{\Cov}{\operatorname{Cov}}
\newcommand{\Var}{\operatorname{Var}}
\newcommand{\tr}{\operatorname{tr}}
\newcommand{\diag}{\operatorname{diag}}
\newcommand{\supp}{\operatorname{supp}}
\newcommand{\law}{\operatorname{law}}
\newcommand{\dd}{\,\mathrm d}
\newcommand{\e}{\mathrm e}
\newcommand{\polylog}{\operatorname{polylog}}
\newcommand{\wtO}{\widetilde O}
\newcommand{\ip}[2]{\left\langle #1,#2\right\rangle}
\newcommand{\norm}[1]{\left\lVert #1\right\rVert}
\newcommand{\abs}[1]{\left|#1\right|}

\title{Heat-capacity annealing for cold-start logconcave sampling}
\author{}
\date{August 2, 2026}

\begin{document}

\maketitle

\begin{abstract}
We study cold-start sampling from a centered isotropic logconcave law
$\pi(\dd x)\propto e^{-V(x)}\dd x$ whose finite convex potential $V$ is
given by an evaluation oracle.  Under the standard ground-set
normalization of Kook and Vempala, we construct,
for output order $q$ and a suitable internal order
$Q=\max\{q,\widetilde O(1)\}$, a sample $Y$ satisfying
$\mathcal R_q(\law(Y)\|\pi)\leq\eps$ using
\[
  \widetilde O\!\left(d^{5/2}\sqrt Q+Qd^2\right)
\]
expected evaluation queries.  For fixed or polylogarithmic order this is
$\widetilde O(d^{2.5})$, improving the
$\widetilde O(qd^{2.75})$ finite-order cold-start bound of Kook and
Vempala on the centered isotropic subclass.

The algorithm is organized by the radial heat capacity of Gaussian
annealing,
\[
  \pi_\beta(\dd x)\propto
  e^{-V(x)-\beta\norm x^2/2}\dd x,
  H(\beta)=\Var_{\pi_\beta}\!\left(\norm X^2\right),
  \qquad H_* = \sup_{\beta\geq0}H(\beta).
\]
For a convex body $K$ with $B(x_0,r)\subseteq K$ and
$\pi=\operatorname{Unif}(K)$, the same argument gives the
heat-parameterized membership-query bound
  \[
  \wtO\!\left(
    d+\frac{d^2\sqrt{QH_*}}{r^2}
    +\frac{Qd^2\norm{\Cov(\pi)}_{\rm op}}{r^2}
  \right).
\]
Here the convex-body parameter is
$H_*=\sup_{\beta\geq0}\Var_{\pi_\beta}\norm{X-x_0}^2$.
The polynomial dependence on $Q$ is displayed; the suppressed factors
are logarithmic in the geometric parameters and the target accuracy.

We prove $H_*=O(\mathcal C_d^3d)$ for every centered isotropic
logconcave law, where $\mathcal C_d$ is the worst isotropic logconcave
Poincar\'e constant.  Klartag's established
$\mathcal C_d=O(\log(ed))$ bound already gives
$H_*=\widetilde O(d)$ and the claimed $d^{2.5}$ exponent.  The key
evaluation-oracle bridge is that Kook and Vempala's fixed exponential-lift
truncation retains at least one half of the mass at every Gaussian
precision; conditioning therefore increases the radial heat capacity by
at most a factor of two.  A recent quadratic-form inequality of Letwin
optionally sharpens $H_*$ to $O(d)$, but removes only a hidden
polylogarithmic factor and is not needed for either sampling exponent.
\end{abstract}

\begin{quote}
\small\noindent\textbf{Note.} This paper was fully AI generated by GPT-5
in the Codex harness, from a ChatGPT Canvas handoff of unspecified model
provenance supplied by Yang Liu, and was prompted by Yang Liu.
\end{quote}

\section{Introduction}

Sampling from a logconcave distribution is one of the basic algorithmic
tasks at the interface of convex geometry, probability, and
optimization.  Its membership-oracle form already contains a canonical
problem.  Given a convex body $K\subset\R^d$, a membership oracle for
$K$, and an accuracy parameter $\eps$, one seeks a random point whose
law is close to the uniform distribution $\pi$ on $K$.  The difficulty
is not merely convergence from a favorable initial law.  A useful
algorithm must also construct such a law from the geometric input---a
\emph{cold start}---without paying an extra factor polynomial in the
dimension.

The modern oracle-complexity theory grew out of volume computation.
Dyer, Frieze, and Kannan gave the first randomized polynomial-time
volume algorithm and an underlying near-uniform sampler
\citep{DFK91random}; Applegate and Kannan initiated the corresponding
theory for general logconcave functions \citep{AK91sampling}.  The
conductance method of Lov\'asz and Simonovits
\citep{LS90mixing,LS93random}, the localization framework of Kannan,
Lov\'asz, and Simonovits \citep{KLS95isop,KLS97random}, and the later
geometry of logconcave functions developed by Lov\'asz and Vempala
\citep{LV07geometry} made random-walk sampling a general-purpose
primitive.  Hit-and-Run, introduced by Smith \citep{Smith84HAR}, was
shown to mix rapidly by Lov\'asz \citep{Lovasz99hit}; the analysis of
Lov\'asz and Vempala permits an arbitrary interior starting point and
separates this guarantee from warm-start analyses
\citep{LV06hit}.  These algorithms differ in their input normalization,
initial-law requirement, output distance, and cost per transition, so
their transition bounds and oracle bounds are not interchangeable.

Annealing supplies the usual bridge from a cold input to a warm law.
Lov\'asz and Vempala used simulated annealing for convex bodies and
logconcave functions \citep{LV06simulated,LV06fast}.  Cousins and
Vempala developed Gaussian cooling for volume and Gaussian-volume
  computation \citep{CV15bypass,CV18Gaussian}.  In a different but related
  literature, thermodynamic-speed schedules based on heat capacity go
  back at least to \citet{SalamonEtAl1988}; adaptive variance and
  partition-function schedules also have a long theoretical computer
  science history, notably \citet{SVV09}.  Our claim is not that
  heat-capacity-guided scheduling is itself new.  The issue
here is whether the particular radial heat capacity of Gaussian cooling
can be controlled strongly enough, and coupled to a finite-R\'enyi
membership-oracle sampler, to improve cold-start query complexity.

The proximal-sampling viewpoint of Lee, Shen, and Tian introduced the
restricted Gaussian oracle as a structured primitive
\citep{LST21structured}.  Chen, Chewi, Salim, and Wibisono sharpened the
analysis of the resulting proximal sampler \citep{CCSW22improved}.
Kook, Vempala, and Zhang implemented this diffusion for convex bodies
through the In-and-Out chain and obtained strong R\'enyi guarantees
\citep{KVZ24INO,KVZ26INO}.  From a cold start, Kook and Zhang obtained
$\widetilde O(d^3)$ membership queries for an
$\mathcal R_\infty$ guarantee \citep{KZ25Renyi}.  Kook and Vempala then
reduced the total-variation warm-start generation cost to
$\widetilde O(d^2R^{3/2}\Lambda^{1/4})$ \citep{KV25faster}, where
\[
  R^2=\E_\pi\norm{X-x_0}^2,
  \qquad \Lambda=\norm{\Cov(\pi)}_{\rm op},
\]
and subsequently obtained the balanced finite-order bound
$\widetilde O(qd^2R^{3/2}\Lambda^{1/4})$ \citep{KV26zeroth}.  In
centered isotropic position,
  $R=\sqrt d$ and $\Lambda=1$, so the latter bound is
  $\widetilde O(qd^{2.75})$, or $\widetilde O(d^{2.75})$ for fixed or
  polylogarithmic order.  Their sampler from a warm start already costs
  $\widetilde O(qd^2\Lambda)$; the $qd^{2.75}$ term is the cost of creating
the warm law, not of mixing from it.  A later unified analysis removes
the $\Lambda\vee1$ loss for arbitrary evaluation-oracle logconcave
densities, giving a $\widetilde O(qd^2\Lambda)$ finite-order
sampling-from-a-warm-start cost under its ground-set normalization
\citep{KV26unified}.

Kook and Vempala also located a genuine obstruction to one tempting
acceleration.  A spherical Gaussian tilt of an isotropic logconcave law
need not preserve the operator norm of the covariance: their
  counterexample has a tilted covariance eigenvalue $\Omega(d^{1/3})$
\citep[Proposition~6.3]{KV26zeroth}.  Thus an analysis that charges every
phase to the worst direction cannot in general keep the covariance
scale constant.  This counterexample does not, however, lower-bound
the scalar fluctuation
\[
  \Var_{\pi_\beta}\norm{X-x_0}^2.
\]
The present paper exploits exactly this distinction.

\subsection{Results}

For $\beta\geq0$, set
\begin{equation}\label{eq:path-intro}
  \pi_\beta(\dd x)
  =\frac{1}{Z(\beta)}
    \exp\!\left(-\frac\beta2\norm{x-x_0}^2\right)\pi(\dd x),
  \qquad
  H(\beta)=\Var_{\pi_\beta}\norm{X-x_0}^2.
\end{equation}
Our algorithmic statement first separates this scalar geometry from
the sampler.

\begin{theorem}[Annealing from a heat-capacity bound, informal]
\label{thm:intro-generic}
Let $B(x_0,r)\subseteq K$, let $\pi$ be uniform on $K$, and suppose
$H(\beta)\leq\overline H$ for every $\beta\geq0$.  Suppose numerical
values $x_0,r$ and upper bounds $\overline H$ and
$\overline\Lambda\geq\norm{\Cov(\pi)}_{\rm op}$ are supplied.  Let
$q\geq1$ be the desired output order and let
$Q\geq\max\{q,2\}$ also meet the polylogarithmic order required by the
finite-R\'enyi restart caps.  There is a cold-start membership-oracle
algorithm whose output $Y$ satisfies
$\mathcal R_q(\law(Y)\|\pi)\leq\eps$ and whose expected query complexity
is
\begin{equation}\label{eq:intro-generic-bound}
  \widetilde O\!\left(
      d+\frac{d^2\sqrt{Q\overline H}}{r^2}
      +\frac{Qd^2\overline\Lambda}{r^2}
  \right).
\end{equation}
No additional polynomial dependence on $Q$ is hidden in the annealing
term.  The
$\widetilde O$ notation suppresses logarithms in the displayed
parameters and in $1/\eps$.
\end{theorem}

The path is parameterized by precision rather than variance.  Its
log-partition function $A(\beta)=\log Z(\beta)$ satisfies
$A''(\beta)=H(\beta)/4$.  Consequently, if $\beta>\beta'$, then
\[
  \mathcal R_Q(\pi_\beta\|\pi_{\beta'})
  \leq \frac Q8(\beta-\beta')^2
       \sup_{\theta\in[\beta',\,\beta'+Q(\beta-\beta')]}H(\theta).
\]
An additive mesh of size $\Theta((Q\overline H)^{-1/2})$ therefore
keeps consecutive targets finitely warm.  A restricted-Gaussian phase
at precision $\beta$ costs $\widetilde O(d^2/(\beta r^2))$ membership
queries.  Summing these costs over the additive grid produces a
harmonic series and the first nontrivial term in
\eqref{eq:intro-generic-bound}.  The last term is the existing balanced
warm-start sampler.

The geometric statement proved here controls the entire path, including the
movement of its mean.

\begin{theorem}[Radial heat capacity]
\label{thm:intro-heat}
Let $\pi$ be any centered isotropic logconcave probability measure on
$\R^d$, and define $\pi_\beta$ by \eqref{eq:path-intro} with $x_0=0$.
Let
\[
  \mathcal C_d=\sup\{\cP(\nu):\nu\text{ is centered, isotropic, and
  logconcave on }\R^d\}.
\]
Then
\begin{equation}\label{eq:intro-poincare-heat}
  \sup_{\beta\geq0}H(\beta)\leq C\mathcal C_d^3d
\end{equation}
for a universal constant $C$.  In particular, the established bound
  $\mathcal C_d=O(\log(ed))$ gives
  $\sup_\beta H(\beta)=O(d\log^3(ed))$.

Using the recent quadratic-form Poincar\'e inequality of Letwin in place
of the general Poincar\'e inequality sharpens the conclusion to
\begin{equation}\label{eq:intro-sharp-heat}
  \sup_{\beta\geq0}H(\beta)=O(d).
\end{equation}
\end{theorem}

The distinction between the two conclusions is deliberate.  The first
part of Theorem~\ref{thm:intro-heat}, combined with Klartag's
  $O(\log(ed))$ Poincar\'e bound \citep{Klartag23log}, already proves the
same $\widetilde O(d^{2.5})$ query exponent.  The much newer theorem of
Letwin \citep{Letwin26KLS} removes a geometric polylogarithmic loss but
is not needed for the headline dimension exponent.  Nor does Letwin's
result prove the full KLS conjecture: it gives a dimension-free
  inequality for quadratic forms and an $O(\log^{1/4}(ed))$ KLS
isoperimetric bound, corresponding to an
  $O(\sqrt{\log(ed)})$ worst Poincar\'e bound under the conventions used
here.  We avoid the common notation collision between that paper's
third-moment parameter and the worst Poincar\'e constant.

More quantitatively, in the normalization $r=\lambda=1$ and with
$Q=\widetilde O(1)$, the $\widetilde O(d^{2.5})$ conclusion needs only
$H_*=\widetilde O(d)$.  Before suppressing logarithmic factors, the
annealing term is $d^2\sqrt{QH_*}$; hence the exact condition for it to
be $o(d^{2.75})$ is $QH_*=o(d^{3/2})$.  In particular, any estimate
$H_*=O(d^{3/2-\delta})$ for a fixed $\delta>0$ gives a genuine
polynomial improvement when $Q$ is polylogarithmic.  Through
$H_*=O(\mathcal C_d^3d)$, the corresponding exact condition is
$Q\mathcal C_d^3=o(d^{1/2})$; for example,
$\mathcal C_d=O(d^{1/6-\delta})$ for any fixed $\delta>0$ is sufficient.
Klartag's logarithmic bound is therefore more than sufficient, and
the Letwin-based linear heat bound is an optional polylogarithmic sharpening
rather than a dependency of the result.

Combining the dependency-minimal form
\eqref{eq:intro-poincare-heat} with
Theorem~\ref{thm:intro-generic} gives the main algorithmic corollary.

\begin{corollary}[Centered scalar-covariance cold start]
\label{cor:intro-main}
Let $K\subset\R^d$ be a convex body given by a membership oracle,
suppose $B(x_0,r)\subseteq K$, and assume
$x_0=\E_\pi X$ and $\Cov(\pi)=\lambda I$, where $x_0,r,\lambda$ are
supplied to the algorithm.  For every $q\geq1$ and
$\eps\in(0,1/100)$, choose
$Q\geq\max\{q,2\}$ satisfying the polylogarithmic cap condition of
Theorem~\ref{thm:algorithm-main}.  There is an algorithm with
\[
  \mathcal R_q(\law(Y)\|\pi)\leq\eps
\]
using
\begin{equation}\label{eq:intro-main-bound}
  \widetilde O\!\left(
    \frac{\lambda}{r^2}
    \bigl(d^{5/2}\sqrt Q+Qd^2\bigr)
  \right)
\end{equation}
expected membership queries.  For fixed-order R\'enyi sampling, or for
total-variation error $\eps$ using an internal
$Q=\widetilde O(1)$, the bound is
$\widetilde O(\lambda d^{2.5}/r^2)$.  In particular it is
$\widetilde O(d^{2.5})$ when $r=\lambda=1$.
\end{corollary}

The same heat estimate accelerates the expensive final phase of the
evaluation-oracle construction.  Write
\[
  L_{\pi,g}=\{x\in\R^d:V(x)-\min V\leq10d\}
\]
for the ground set used in that construction.

\begin{theorem}[Centered isotropic evaluation-oracle cold start]
\label{thm:intro-evaluation}
Let $V:\R^d\to\R$ be a convex function given by an evaluation oracle.
Suppose
$\pi(\dd x)\propto e^{-V(x)}\dd x$ is centered and isotropic and
$B(0,1)\subseteq L_{\pi,g}$.  For every $q\geq1$ and
$\eps\in(0,1/100)$, there is a cold-start algorithm returning $Y$ with
\[
  \mathcal R_q(\law(Y)\|\pi)\leq\eps
\]
after
\begin{equation}\label{eq:intro-evaluation-bound}
  \widetilde O\!\left(d^{5/2}\sqrt Q+Qd^2\right)
\end{equation}
expected evaluation queries, where
$Q=\max\{q,\widetilde O(1)\}$ meets the finite-order cap conditions.
In particular, the complexity is $\widetilde O(d^{2.5})$ for fixed or
polylogarithmic output order.
\end{theorem}

Theorem~\ref{thm:intro-evaluation} uses the same fixed truncation as
Kook and Vempala; it does not replace their exponential lift or their
sampler.  The additional ingredient in our argument is a uniform
comparison between the truncated and untruncated Gaussian paths,
described below and proved in Section~\ref{sec:lift}.

The centering and scalar-covariance hypotheses are substantive.  A
spherical quadratic tilt is not affine invariant, so whitening does not
turn an arbitrary anisotropic instance into Theorem~\ref{thm:intro-heat}.
We do not claim the scale-aware bound
$H_*=O(R^2\Lambda)$ for general covariance.

\subsection{Proof overview}

Write $m_\beta$ and $\Sigma_\beta$ for the mean and covariance of
$\pi_\beta$, and introduce
\begin{equation}\label{eq:intro-potentials}
  G_\beta=\tr(\Sigma_\beta^2)
     +m_\beta^\top\Sigma_\beta m_\beta,
  \qquad
  \Phi_\beta=-\log\det\Sigma_\beta.
\end{equation}
Energy monotonicity gives
$\tr\Sigma_\beta+\norm{m_\beta}^2\leq d$, while Brascamp--Lieb gives
$\Sigma_\beta\preceq\beta^{-1}I$.  A deterministic eigenvalue
calculation then yields
\begin{equation}\label{eq:intro-G-Phi}
  G_\beta\leq d+\frac{3\Phi_\beta}{\beta}.
\end{equation}
The exact covariance derivative gives
\[
  \Phi_\beta'
  =\frac12\Cov_{\pi_\beta}\!\left(
      (X-m_\beta)^\top\Sigma_\beta^{-1}(X-m_\beta),
      \norm X^2
    \right).
\]
After whitening $\pi_\beta$, a general Poincar\'e inequality bounds the
two variances in this covariance by
$4\mathcal C_dd$ and $4\mathcal C_dG_\beta$.  Hence
$\abs{\Phi_\beta'}\leq2\mathcal C_d\sqrt{dG_\beta}$.  A first-crossing
argument against a linear barrier gives
$\Phi_\beta=O(\mathcal C_d^2d\beta)$, then
$G_\beta=O(\mathcal C_d^2d)$ and
$H(\beta)=O(\mathcal C_d^3d)$.

Letwin's quadratic-form theorem can be inserted at exactly two points.
It gives $H(\beta)\leq16G_\beta$ and bounds the whitened Mahalanobis
variance by $8d$.  The same first-crossing argument then gives
$\Phi_\beta=O(d\beta)$ and $H(\beta)=O(d)$.  This route is cleaner, but
the preceding Poincar\'e route is retained to show that the sampling
exponent is robust to substantially weaker geometry.

\subsection{General evaluation-oracle densities}

For $\pi(\dd x)\propto e^{-V(x)}\dd x$, Kook and Vempala's exponential
lift at Gaussian precision $\beta$ has density
\[
  P_\beta(\dd x\dd t)\propto
  e^{-\beta\norm x^2/2-dt}\1_{\{V(x)\leq dt\}}\dd x\dd t.
\]
Its $x$-marginal is $\pi_\beta$, and conditional on $x$ the $t$-law is
independent of $\beta$.  The cold-start construction restricts the
epigraph to one fixed bounded convex cylinder $A$.  Let
$p_\beta=P_\beta(A)$ and let $\nu_\beta$ be the $x$-marginal of
$P_\beta(\cdot\mid A)$.  The decisive observation is
\begin{equation}\label{eq:intro-uniform-truncation}
  p_\beta\geq\frac12\quad(\beta\geq0),
  \qquad
  \Var_{\nu_\beta}\norm X^2
  \leq2\Var_{\pi_\beta}\norm X^2.
\end{equation}
To see the first inequality, put
$W_\beta(x)=V(x)+\beta\norm x^2/2$ and
$s=t+\beta\norm x^2/(2d)$.  The triangular change of variables
$(x,t)\mapsto(x,s)$ sends $P_\beta$ to the standard exponential lift of
$W_\beta$.  Kook and Vempala's same convex-truncation lemma applies to
$W_\beta$: energy monotonicity decreases its radial second moment,
$W_\beta(0)=V(0)$ leaves the upper lift cutoff unchanged, and $t\leq s$.
Thus a set of mass at least one half for the transformed lift lies
inside $A$.  The second inequality follows from
$\dd\nu_\beta/\dd\pi_\beta\leq p_\beta^{-1}\leq2$ by centering the
squared radius at its $\pi_\beta$ mean.

Consequently the Poincar\'e heat bound for the exact isotropic target
controls the whole implemented truncated path.  Replacing only the
third phase of Kook and Vempala's warm-start construction by the
additive precision schedule costs
$\widetilde O(d^2\sqrt{QH_*})$.  At precision zero the truncation is
still $2$-warm for the full lift, so their balanced full-lift sampler
finishes the algorithm in $\widetilde O(Qd^2)$ queries.

\paragraph{Organization.}
Section~\ref{sec:preliminaries} fixes conventions and records the
finite-R\'enyi sampler inputs.  Section~\ref{sec:annealing} proves the
exact overlap identity and the heat-capacity-parameterized algorithm.
Section~\ref{sec:heat} proves the Poincar\'e and quadratic-form heat
bounds.  Section~\ref{sec:consequences} gives the cold-start corollaries
and comparisons.  Section~\ref{sec:lift} proves the uniform-truncation
lemma and the evaluation-oracle theorem.

\section{Preliminaries and sampler inputs}
\label{sec:preliminaries}

All probability measures are full-dimensional unless explicitly
stated otherwise.  A probability measure on $\R^d$ is logconcave if it
has a density $e^{-V}$ on its affine support for a convex extended-real
function $V$.  A logconcave law is \emph{isotropic} if it has mean zero
and covariance $I$.

For a probability measure $\nu$ with finite Poincar\'e constant, write
$\cP(\nu)$ for the least $C$ such that
\begin{equation}\label{eq:poincare}
  \Var_\nu f\leq C\E_\nu\norm{\nabla f}^2
\end{equation}
for every locally Lipschitz $f\in L^2(\nu)$.  Testing linear functions
shows that the worst isotropic constant $\mathcal C_d$ defined in
Theorem~\ref{thm:intro-heat} is at least one.  The KLS conjecture asks
whether $\sup_d\mathcal C_d<\infty$.  We only use the proved bound
$\mathcal C_d=O(\log(ed))$ from Klartag's isoperimetric theorem and the
standard equivalence between the Cheeger and Poincar\'e constants for
logconcave measures \citep{Chee70lower,Buser82iso,Led04spectral,Klartag23log}.

For $a>1$ and probability measures $\mu\ll\nu$, the order-$a$ R\'enyi
divergence is
\begin{equation}\label{eq:renyi}
  \mathcal R_a(\mu\|\nu)
  =\frac{1}{a-1}\log\int
    \left(\frac{\dd\mu}{\dd\nu}\right)^a\dd\nu.
\end{equation}
We use the continuous extension $\mathcal R_1(\mu\|\nu)
=\operatorname{KL}(\mu\|\nu)$ at order one.
It is nondecreasing in $a$ and obeys data processing.  The associated
warmness is
\[
  M_a(\mu\|\nu)
  =\norm{\dd\mu/\dd\nu}_{L^a(\nu)}
  =\exp\!\left(\frac{a-1}{a}\mathcal R_a(\mu\|\nu)\right).
\]
We use the following weak triangle inequality
\citep{vH14renyi,KV26zeroth}.

\begin{lemma}[Weak R\'enyi triangle inequality]
\label{lem:weak-triangle}
For $a>1$, $\theta\in(0,1)$, and probability measures
$\mu,\nu,\rho$,
\begin{equation}\label{eq:weak-triangle}
  \mathcal R_a(\mu\|\rho)
  \leq
  \frac{a-\theta}{a-1}
    \mathcal R_{a/\theta}(\mu\|\nu)
  +\mathcal R_{(a-\theta)/(1-\theta)}(\nu\|\rho).
\end{equation}
In particular, with $\theta=a/(2a-1)$, both orders on the right are at
most $2a$, and the first coefficient is at most two.
\end{lemma}

The choice of $\theta$ in the last sentence is important.  It is the
order-dependent choice that justifies relaying a $2a$ guarantee through
arbitrarily many phases; a fixed choice such as $2/3$ does not give the
displayed pair of orders for every $a$.

We next state only the sampler consequences used later.  They follow
from the capped restricted-Gaussian sampler and the balanced uniform
sampler of Kook and Vempala
\citep[Theorems~1.1 and~1.4 and Proposition~4.1]{KV26zeroth}.  We
distinguish the ideal uncapped chain, the capped implementation, and the
restart wrapper, since conditioning these laws interchangeably would
lose the final accuracy.  The exact logarithmic factors and rejection
thresholds are retained in the cited source.

\begin{theorem}[Finite-order sampler inputs]
\label{thm:sampler-inputs}
Let $B(x_0,r)\subseteq K$ and let
\[
  \mu_\beta(\dd x)\propto
  \1_K(x)e^{-\beta\norm{x-x_0}^2/2}\dd x,
  \qquad 0<\beta\leq d/r^2.
\]
Fix target error $\zeta\in(0,1/10)$ and cap probability
$\eta\in(0,1/2)$.  There are universal constants with the following
properties.

\begin{enumerate}[label=(\roman*)]
\item Suppose an initial law has $M_Q\leq10$ with respect to
  $\mu_\beta$, where $Q\geq2$.  After a chosen horizon $N_G$, the ideal
  uncapped restricted-Gaussian chain has $\mathcal R_{2Q}$ error at most
  $\zeta$.  Its capped implementation agrees with that chain unless a
  cap is hit, has cap-hit probability at most $\eta$, and has
  unconditional expected membership-query cost
  \begin{equation}\label{eq:gaussian-phase-cost}
    \widetilde O\!\left(\frac{d^2}{\beta r^2}\right).
  \end{equation}
  It is enough that $Q\geq C\log(N_G/\eta)$.  No polynomial factor in
  $Q$ is suppressed in \eqref{eq:gaussian-phase-cost}.

\item Let $\pi$ be uniform on $K$ and
  $\Lambda=\norm{\Cov(\pi)}_{\rm op}$.  From an initial law with
  $M_Q\leq10$, the ideal uncapped balanced uniform chain has
  $\mathcal R_Q$ error at most $\zeta$ after a chosen horizon $N_U$.
  The capped implementation agrees with it off an event of probability
  at most $\eta$ and has unconditional expected cost
  \begin{equation}\label{eq:uniform-phase-cost}
    \widetilde O\!\left(
       \frac{Qd^2\Lambda}{r^2}\right)
  \end{equation}
  membership queries, provided $Q\geq C\log(N_U/\eta)$.  Here the
  native bound contains a factor $\log^6(1/\zeta)$; the remaining
  suppressed factors are logarithmic.
\end{enumerate}

The horizons are the actual iteration counts selected by the algorithms;
each is bounded by its displayed query bound.  A restart wrapper must
draw a fresh input on every attempt.  If $S$ is the event that a coupled
capped run never hits a cap, its returned law is the ideal output
conditioned on $S$.
\end{theorem}

After the rescaling $y=(x-x_0)/r$, the variance of the restricted
Gaussian is $1/(\beta r^2)$, which explains
\eqref{eq:gaussian-phase-cost}.  The range
$\beta\leq d/r^2$ ensures that this variance is at least $1/d$, the
initial scale used below.

\section{Annealing controlled by radial heat capacity}
\label{sec:annealing}

Let $\pi$ be any probability measure for which the following
normalizing constants are finite.  Put
\[
  T(x)=\norm{x-x_0}^2,
  \qquad
  Z(\beta)=\int e^{-\beta T(x)/2}\pi(\dd x),
  \qquad
  A(\beta)=\log Z(\beta),
\]
and let $\pi_\beta$ be the normalized tilted law.  For a uniform law on
a compact body, these quantities are analytic for every
$\beta\geq0$.  For a general logconcave law, the same identities follow
by truncation and dominated convergence; logconcavity gives all
polynomial moments.

\begin{lemma}[Heat-capacity identities]
\label{lem:partition-identities}
For every $\beta\geq0$ at which the derivatives are taken,
\begin{equation}\label{eq:A-derivatives}
  A'(\beta)=-\frac12\E_{\pi_\beta}T,
  \qquad
  A''(\beta)=\frac14\Var_{\pi_\beta}T=\frac14H(\beta).
\end{equation}
If $\beta\geq\beta'\geq0$ and $a>1$, then
\begin{equation}\label{eq:exact-renyi}
  \mathcal R_a(\pi_\beta\|\pi_{\beta'})
  =\frac{
    A(\beta'+a(\beta-\beta'))+(a-1)A(\beta')-aA(\beta)
  }{a-1}.
\end{equation}
Consequently,
\begin{equation}\label{eq:renyi-heat-bound}
  \mathcal R_a(\pi_\beta\|\pi_{\beta'})
  \leq\frac a8(\beta-\beta')^2
  \sup_{\theta\in[\beta',\,\beta'+a(\beta-\beta')]}H(\theta).
\end{equation}
The same statement with $\beta'=0$ compares a Gaussian tilt with the
untilted target.
\end{lemma}

\begin{proof}
Differentiating under the integral gives
\eqref{eq:A-derivatives}.  If $p_\gamma$ denotes the density of
$\pi_\gamma$ with respect to $\pi$, then
\begin{align*}
 \int p_\beta^a p_{\beta'}^{1-a}\dd\pi
 &=\exp\!\left(
 A(\beta'+a(\beta-\beta'))-aA(\beta)+(a-1)A(\beta')
 \right),
\end{align*}
which proves \eqref{eq:exact-renyi}.  The numerator in
\eqref{eq:exact-renyi} vanishes for affine $A$.  Applying the integral
remainder formula to its second derivative, and using
$A''=H/4$, bounds it by
$a(a-1)(\beta-\beta')^2\sup H/8$.  Division by $a-1$ proves
\eqref{eq:renyi-heat-bound}.
\end{proof}

The direction in \eqref{eq:exact-renyi} is the one needed by cooling:
the old, more concentrated law is measured relative to the new, less
concentrated target.  Reversing the divergence probes the parameter
$a\beta'-(a-1)\beta$, which may be negative and need not even define a
normalizable tilt for an unbounded base law.

\subsection{The additive schedule}

We now prove the formal version of
Theorem~\ref{thm:intro-generic}.  The theorem is stated with an internal
order $Q$ because small requested orders must be raised to meet the
logarithmic rejection-cap threshold.  This is a polylogarithmic issue,
not a hidden polynomial factor.

\begin{theorem}[Heat-capacity annealing]
\label{thm:algorithm-main}
Let $K\subset\R^d$ be a convex body, $B(x_0,r)\subseteq K$, and let
$\pi$ be uniform on $K$.  Suppose the algorithm is supplied with valid
numerical values $x_0,r$ and bounds $\overline H$ and
$\overline\Lambda$ such that
\begin{equation}\label{eq:Hbar-assumption}
  H(\beta)=\Var_{\pi_\beta}\norm{X-x_0}^2\leq\overline H
  \quad(\beta\geq0),
  \qquad
  \norm{\Cov(\pi)}_{\rm op}\leq\overline\Lambda.
\end{equation}
Fix $q\geq1$ and $\eps\in(0,1/100)$.  Define the dimensionless quantity
\[
  \Xi=e^2+d+q+\frac{\sqrt{\overline H}}{r^2}
        +\frac{\overline\Lambda}{r^2}+\frac1\eps.
\]
For a sufficiently large universal constant $C_0$, take
\begin{equation}\label{eq:cap-conditions}
  Q\geq\max\{q,2,C_0\log^2\Xi\}.
\end{equation}
This explicit choice dominates the logarithmic iteration-and-cap
thresholds in Theorem~\ref{thm:sampler-inputs}; in particular,
$Q=\max\{q,\widetilde O(1)\}$.  There is a restart-until-success
algorithm whose output $Y$ satisfies
\[
  \mathcal R_q(\law(Y)\|\pi)\leq\eps
\]
and whose expected membership-query complexity is
\begin{equation}\label{eq:generic-query-bound}
  \widetilde O\!\left(
      d+\frac{d^2\sqrt{Q\overline H}}{r^2}
      +\frac{Qd^2\overline\Lambda}{r^2}
  \right).
\end{equation}
The algorithm terminates almost surely.  A finite-attempt version may
instead declare failure with any prescribed probability $\delta>0$,
at the cost of an additional logarithmic factor in $1/\delta$.
\end{theorem}

\begin{proof}
We first describe an ideal uncapped pipeline and then couple all caps to
it.  Set
\[
  \beta_0=\frac d{r^2},
  \qquad
  \Delta=\min\!\left\{\beta_0,
    \frac{1}{4\sqrt{Q\overline H}}\right\},
  \qquad
  L=\left\lceil\frac{\beta_0}{\Delta}\right\rceil,
  \qquad
  h=\frac{\beta_0}{L}.
\]
Because $\Delta\leq\beta_0$, we have
\begin{equation}\label{eq:grid-comparison}
  \frac{\Delta}{2}\leq h\leq\Delta,
  \qquad
  L\leq2+\frac{4d\sqrt{Q\overline H}}{r^2}.
\end{equation}
Use the positive levels
\[
  \beta_i=(L-i)h,
  \qquad i=0,\ldots,L-1,
\]
followed by the untilted target at zero.  Thus the smallest positive
precision is $h$, rather than an uncontrolled residual below the mesh.

Draw exactly from $\pi_{\beta_0}$ by rejection sampling from
$N(x_0,r^2I/d)$.  The acceptance probability is at least
$\Pp(\chi_d^2\leq d)>1/2$, because
$B(x_0,r)\subseteq K$; hence the expected initialization cost is
universal.

For consecutive exact targets, Lemma~\ref{lem:partition-identities} at
order $2Q$ gives
\begin{equation}\label{eq:neighbor-warm}
  \mathcal R_{2Q}(\pi_{\beta_i}\|\pi_{\beta_{i+1}})
  \leq \frac{2Q}{8}h^2\overline H
  \leq\frac1{64}.
\end{equation}
Since $\beta_{L-1}=h\leq\Delta$, the same calculation with
$\beta'=0$ gives
\begin{equation}\label{eq:final-warm}
  \mathcal R_{2Q}(\pi_{\beta_{L-1}}\|\pi)
  \leq\frac1{64}.
\end{equation}

Let $\widehat\pi_0=\pi_{\beta_0}$.  Inductively, run the ideal uncapped
restricted-Gaussian chain toward $\pi_{\beta_{i+1}}$ long enough that
its output $\widehat\pi_{i+1}$ satisfies
$\mathcal R_{2Q}(\widehat\pi_{i+1}\|\pi_{\beta_{i+1}})\leq c_0$, where
$c_0$ is a sufficiently small universal constant.  Indeed, if the
invariant holds at level $i$, Lemma~\ref{lem:weak-triangle} with
$a=Q$ and $\theta=Q/(2Q-1)$, monotonicity, and
\eqref{eq:neighbor-warm} give
\[
  \mathcal R_Q(\widehat\pi_i\|\pi_{\beta_{i+1}})
  \leq 2c_0+\frac1{64}<1.
\]
Thus the next call starts with $M_Q<e<10$, and
Theorem~\ref{thm:sampler-inputs}(i) restores the order-$2Q$ invariant.
Equation~\eqref{eq:final-warm} gives the same constant warmness with
respect to $\pi$.  Starting from $\widehat\pi_{L-1}$, run the ideal
uncapped balanced uniform chain to an output law $\rho$ satisfying
\begin{equation}\label{eq:ideal-final-error}
  \mathcal R_Q(\rho\|\pi)\leq\frac\eps2.
\end{equation}

We next sum the unconditional capped-call costs that will bound one
full attempt.  Reversing the positive grid, the Gaussian phases have
precisions $kh$, $k=1,\ldots,L-1$.  Therefore
\begin{align}
 \sum_{k=1}^{L-1}\widetilde O\!\left(
      \frac{d^2}{khr^2}\right)
 &\leq
 \widetilde O\!\left(\frac{d^2}{hr^2}\right)\notag\\
 &\leq\widetilde O\!\left(
   d+\frac{d^2\sqrt{Q\overline H}}{r^2}
 \right).
 \label{eq:harmonic-cost}
\end{align}
Here we used \eqref{eq:grid-comparison} and
$1/\Delta=\max\{r^2/d,4\sqrt{Q\overline H}\}$.  The uniform phase adds
$\widetilde O(Qd^2\overline\Lambda/r^2)$.

Finally, put caps back into the coupled pipeline.  Set
$\eta=\eps/8$, give every one of the at most $L-1$ Gaussian phases cap
budget $\eta/(2L)$, and give the uniform phase budget $\eta/2$.  The
actual horizons $N_G$ and $N_U$ are at most the corresponding displayed
query bounds in Theorem~\ref{thm:sampler-inputs}.  Using
\eqref{eq:grid-comparison}, their logarithms, together with
$\log(L/\eta)$, are bounded by a universal constant times
$\log\Xi+\log Q$.  Thus the explicit choice
\eqref{eq:cap-conditions}, for sufficiently large $C_0$, meets every
condition $Q\geq C\log(N/\eta')$ in that theorem.  A union bound shows
that the event $S$ of no cap hit in a full attempt has probability at
least $1-\eta$.

On a cap hit, discard the attempt---including any completed annealing
phases---and restart from a fresh exact sample at $\beta_0$.  This also
  supplies the fresh warm input required when the final uniform chain is
  restarted.  The unconditional cost of one capped attempt is at most
  the sum of the displayed capped-call costs, since a phase is run only
  when every preceding cap has succeeded.  The expected number of
  attempts is at most $(1-\eta)^{-1}$.

It remains to check the law, rather than only the cost.  Let $\nu$ be
the output law of a successful attempt.  The capped and ideal pipelines
agree on $S$, so
\[
  \frac{\dd\nu}{\dd\rho}\leq\frac1{\Pp(S)}
  \leq\frac1{1-\eta}.
\]
Directly from the definition of R\'enyi divergence and
\eqref{eq:ideal-final-error},
\[
  \mathcal R_Q(\nu\|\pi)
  \leq\frac{\eps}{2}+\frac{Q}{Q-1}\log\frac1{1-\eta}
  \leq\eps,
\]
where the last inequality uses $Q\geq2$, $\eta=\eps/8$, and
$\eps<1/100$.  Independent full attempts therefore terminate almost
surely and return exactly the law $\nu$; monotonicity gives the desired
order-$q$ statement.  Stopping after a logarithmic number of attempts
gives the asserted finite-attempt variant.
\end{proof}

\begin{remark}[Total variation]
By Pinsker's inequality,
$2\TV(\mu,\pi)^2\leq\mathcal R_1(\mu\|\pi)\leq\mathcal R_Q(\mu\|\pi)$.
To obtain total-variation error $\eps$, run
Theorem~\ref{thm:algorithm-main} with R\'enyi error $2\eps^2$.  This
changes only logarithmic accuracy factors.
\end{remark}

\section{Uniform radial heat capacity}
\label{sec:heat}

This section proves Theorem~\ref{thm:intro-heat}.  We work first with a
centered isotropic logconcave probability measure $\pi$ on $\R^d$ and
write
\[
  \pi_\beta(\dd x)\propto e^{-\beta\norm x^2/2}\pi(\dd x).
\]
Let $m=m_\beta$ and $\Sigma=\Sigma_\beta$ denote its mean and
covariance.  Dependence on $\beta$ is suppressed locally.  Define
\begin{equation}\label{eq:G-Phi}
  G_\beta=\tr(\Sigma_\beta^2)
     +m_\beta^\top\Sigma_\beta m_\beta,
  \qquad
  \Phi_\beta=-\log\det\Sigma_\beta.
\end{equation}
Logconcavity gives all polynomial moments.  Consequently the moment
maps used below are right-differentiable at $\beta=0$ and continuously
differentiable for $\beta>0$; in particular, the first-contact arguments
below apply to $\Phi_\beta$.

\subsection{Energy, determinant, and covariance identities}

\begin{lemma}[Energy monotonicity]
\label{lem:energy}
For every $\beta\geq0$,
\begin{equation}\label{eq:energy-monotonicity}
  \frac{\dd}{\dd\beta}\E_{\pi_\beta}\norm X^2
  =-\frac12H(\beta),
  \qquad
  \tr\Sigma_\beta+\norm{m_\beta}^2\leq d.
\end{equation}
\end{lemma}

\begin{proof}
The derivative is the standard exponential-family covariance identity,
or follows from Lemma~\ref{lem:partition-identities}.  At $\beta=0$,
centered isotropy gives $\E_\pi\norm X^2=d$; the derivative is
nonpositive.
\end{proof}

\begin{lemma}[Brascamp--Lieb and the deterministic potential bound]
\label{lem:deterministic-G}
For $\beta>0$,
\begin{equation}\label{eq:BL-covariance}
  \Sigma_\beta\preceq\beta^{-1}I,
\end{equation}
and
\begin{equation}\label{eq:G-bound}
  \Phi_\beta\geq d-\tr\Sigma_\beta\geq\norm{m_\beta}^2,
  \qquad
  G_\beta\leq d+\frac{3\Phi_\beta}{\beta}.
\end{equation}
\end{lemma}

\begin{proof}
The potential defining $\pi_\beta$ is $\beta$-strongly convex, allowing
an extended-real convex part.  The Brascamp--Lieb variance inequality
therefore gives \eqref{eq:BL-covariance}
\citep{BrascampLieb76}; an indicator of a convex support follows by
standard approximation.

Let $\lambda_1,\ldots,\lambda_d$ be the eigenvalues of $\Sigma_\beta$
and set $h(t)=t-1-\log t$.  Since $h\geq0$,
\[
  \sum_i h(\lambda_i)
  =\tr\Sigma_\beta-d+\Phi_\beta\geq0.
\]
Together with Lemma~\ref{lem:energy}, this proves
$\Phi_\beta\geq d-\tr\Sigma_\beta\geq\norm{m_\beta}^2$ and also
$\sum_i h(\lambda_i)\leq\Phi_\beta$.

For every $t>0$,
\begin{equation}\label{eq:scalar-entropy}
  (t-1)^2\leq2\max\{1,t\}\,h(t).
\end{equation}
If $0<\beta\leq1$, then
$\max\{1,\lambda_i\}\leq1/\beta$, and
\begin{align*}
  \tr(\Sigma_\beta^2)
  &=d+2(\tr\Sigma_\beta-d)+\sum_i(\lambda_i-1)^2\\
  &\leq d+\frac{2\Phi_\beta}{\beta}.
\end{align*}
If $\beta\geq1$, then
$\tr(\Sigma_\beta^2)\leq\beta^{-1}\tr\Sigma_\beta\leq d$, so the
same displayed upper bound remains true.  Finally,
\[
  m_\beta^\top\Sigma_\beta m_\beta
  \leq\frac{\norm{m_\beta}^2}{\beta}
  \leq\frac{\Phi_\beta}{\beta}.
\]
Adding the two estimates proves \eqref{eq:G-bound}.
\end{proof}

\begin{lemma}[Determinant derivative]
\label{lem:Phi-derivative}
For $\beta\geq0$, with the identity at zero interpreted as a right
derivative,
\begin{equation}\label{eq:Sigma-derivative}
  \Sigma_\beta'
  =-\frac12\Cov_{\pi_\beta}\!\left(
     (X-m_\beta)(X-m_\beta)^\top,\norm X^2
  \right),
\end{equation}
and
\begin{equation}\label{eq:Phi-derivative}
  \Phi_\beta'
  =\frac12\Cov_{\pi_\beta}\!\left(S_\beta,T\right),
  \quad
  S_\beta=(X-m_\beta)^\top\Sigma_\beta^{-1}(X-m_\beta),
  \quad T=\norm X^2.
\end{equation}
\end{lemma}

\begin{proof}
For an integrable statistic $F$ that does not depend on $\beta$,
\[
  \frac{\dd}{\dd\beta}\E_{\pi_\beta}F
  =-\frac12\Cov_{\pi_\beta}(F,T).
\]
Differentiate the centered covariance.  The two terms produced by
$m_\beta'$ vanish because
$\E_{\pi_\beta}(X-m_\beta)=0$, which proves
\eqref{eq:Sigma-derivative}.  The identity
$\Phi_\beta'=-\tr(\Sigma_\beta^{-1}\Sigma_\beta')$ then gives
\eqref{eq:Phi-derivative}.
\end{proof}

No monotonicity of $\Phi_\beta$ is asserted or needed.  The argument
below controls its absolute derivative.

\subsection{A Poincar\'e proof sufficient for the query exponent}

\begin{theorem}[Poincar\'e heat bound]
\label{thm:poincare-heat}
For every centered isotropic logconcave $\pi$ on $\R^d$ and every
$\beta\geq0$,
\begin{equation}\label{eq:poincare-heat-result}
  \Phi_\beta\leq16\mathcal C_d^2d\beta,
  \qquad
  G_\beta\leq49\mathcal C_d^2d,
  \qquad
  H(\beta)\leq196\mathcal C_d^3d.
\end{equation}
\end{theorem}

\begin{proof}
Let
$Y=\Sigma_\beta^{-1/2}(X-m_\beta)$.  Its law is isotropic and
logconcave.  As a function of $Y$,
\[
  T=\norm{m_\beta+\Sigma_\beta^{1/2}Y}^2,
  \qquad
  \E\norm{\nabla_YT}^2
  =4\left(\tr(\Sigma_\beta^2)
      +m_\beta^\top\Sigma_\beta m_\beta\right)
  =4G_\beta.
\]
The Poincar\'e inequality therefore gives
\begin{equation}\label{eq:H-via-G-PI}
  H(\beta)\leq4\mathcal C_dG_\beta.
\end{equation}
Likewise, $S_\beta=\norm Y^2$, so
\begin{equation}\label{eq:S-variance-PI}
  \Var(S_\beta)\leq4\mathcal C_dd.
\end{equation}
By Cauchy--Schwarz in \eqref{eq:Phi-derivative},
\begin{equation}\label{eq:Phi-diff-PI}
  \abs{\Phi_\beta'}
  \leq2\mathcal C_d\sqrt{dG_\beta}.
\end{equation}

  We claim that
  $\Phi_\beta\leq16\mathcal C_d^2d\beta$.  At zero,
  $\Phi_0=0$, $G_0=d$, and \eqref{eq:H-via-G-PI} gives
  $H(0)\leq4\mathcal C_dd$.  Moreover, $S_0=T=\norm X^2$ because
  $m_0=0$ and $\Sigma_0=I$, so
  $\Phi_0'=H(0)/2\leq2\mathcal C_dd$, which is strictly below the
proposed barrier slope because $\mathcal C_d\geq1$.  If there were a
first positive contact with the barrier, then
Lemma~\ref{lem:deterministic-G} and \eqref{eq:Phi-diff-PI} would give
\[
  \Phi_\beta'
  \leq2\mathcal C_dd\sqrt{1+48\mathcal C_d^2}
  \leq14\mathcal C_d^2d,
\]
whereas first contact from below requires
$\Phi_\beta'\geq16\mathcal C_d^2d$.  This contradiction proves the
first estimate.  Substitution into \eqref{eq:G-bound} gives
$G_\beta\leq49\mathcal C_d^2d$, and
\eqref{eq:H-via-G-PI} gives the last estimate.
\end{proof}

Combining Theorem~\ref{thm:poincare-heat} with
$\mathcal C_d=O(\log(ed))$ proves
$H_*=O(d\log^3(ed))$.  This is already enough for
$\widetilde O(d^{2.5})$ sampling.

The following subsection improves the heat bound itself, but is not
needed for either sampling exponent.

\subsection{Optional sharpening from quadratic-form variance}

The input is the following recent preprint theorem.  We state its
status because the preceding subsection deliberately avoids depending
on it for the query exponent.

\begin{theorem}[Letwin's quadratic-form inequality]
\label{thm:letwin}
The following is Theorem~1.2 of \citet{Letwin26KLS}.
If $Y$ is isotropic and logconcave on $\R^d$, then for every symmetric
matrix $M$,
\begin{equation}\label{eq:letwin}
  \Var(Y^\top MY)
  \leq2\E\norm{\nabla(Y^\top MY)}^2
  =8\tr(M^2).
\end{equation}
The theorem applies without smoothness or bounded-support assumptions,
and without an invertibility assumption on $M$.
\end{theorem}

\begin{theorem}[Optimal-order radial heat bound]
\label{thm:sharp-heat}
Under the hypotheses of Theorem~\ref{thm:poincare-heat},
\begin{equation}\label{eq:sharp-heat-result}
  \Phi_\beta\leq128d\beta,
  \qquad
  G_\beta\leq385d,
  \qquad
  H(\beta)\leq6160d
  \qquad(\beta\geq0).
\end{equation}
\end{theorem}

\begin{proof}
Again let $Y=\Sigma_\beta^{-1/2}(X-m_\beta)$.  Decompose
\begin{equation}\label{eq:T-decomposition}
  T=\norm{m_\beta}^2
    +Y^\top\Sigma_\beta Y
    +2(\Sigma_\beta^{1/2}m_\beta)^\top Y.
\end{equation}
Theorem~\ref{thm:letwin}, the exact variance of the linear term, and
$\Var(U+V)\leq2\Var U+2\Var V$ yield
\begin{align}
  H(\beta)
  &\leq16\tr(\Sigma_\beta^2)
       +8m_\beta^\top\Sigma_\beta m_\beta
   \leq16G_\beta.
  \label{eq:H-via-G-sharp}
\end{align}
Applying Theorem~\ref{thm:letwin} to $S_\beta=Y^\top IY$ gives
$\Var(S_\beta)\leq8d$.  Hence
\begin{equation}\label{eq:Phi-diff-sharp}
  \abs{\Phi_\beta'}
  \leq\frac12\sqrt{8d\cdot16G_\beta}
  =4\sqrt{2dG_\beta}.
\end{equation}

Consider the barrier $128d\beta$.  At zero, $G_0=d$ and
Theorem~\ref{thm:letwin} gives $H(0)\leq8d$.  As above, $S_0=T$, so
$\Phi_0'=H(0)/2\leq4d$.  At a hypothetical first positive contact,
Lemma~\ref{lem:deterministic-G} would give $G_\beta\leq385d$, and
\eqref{eq:Phi-diff-sharp} would give
\[
  \Phi_\beta'\leq4\sqrt{770}\,d<112d<128d,
\]
contradicting first contact.  Thus $\Phi_\beta\leq128d\beta$.
The remaining two estimates follow from \eqref{eq:G-bound} and
\eqref{eq:H-via-G-sharp}.
\end{proof}

The order $d$ cannot be improved.  If $\pi$ is uniform on
$[-\sqrt3,\sqrt3]^d$, then it is isotropic and
\[
  H(0)=d\Var(U^2)=\frac45d,
  \qquad U\sim\operatorname{Unif}[-\sqrt3,\sqrt3].
\]

\begin{corollary}[Scalar covariance]
\label{cor:scalar-heat}
Let $\pi$ be logconcave with mean $x_0$ and
$\Cov(\pi)=\lambda I$.  Along the spherical path
\eqref{eq:path-intro},
\begin{equation}\label{eq:scalar-heat}
  \sup_{\beta\geq0}
  \Var_{\pi_\beta}\norm{X-x_0}^2
  \leq C d\lambda^2.
\end{equation}
Without Theorem~\ref{thm:letwin}, the right side may be replaced by
$C\mathcal C_d^3d\lambda^2$.
\end{corollary}

\begin{proof}
Translate by $x_0$ and scale by $\lambda^{-1/2}$.  The transformed base
law is isotropic, the precision changes from $\beta$ to
$\lambda\beta$, and the squared-radius statistic changes by the factor
$\lambda$.  Its variance therefore changes by $\lambda^2$.
\end{proof}

\section{Cold-start consequences and comparisons}
\label{sec:consequences}

\begin{theorem}[Cold-start sampling with scalar covariance]
\label{thm:cold-main}
Let $K\subset\R^d$ be given by a membership oracle and let
$\pi=\operatorname{Unif}(K)$.  Suppose
\[
  B(x_0,r)\subseteq K,
  \qquad x_0=\E_\pi X,
  \qquad \Cov(\pi)=\lambda I.
\]
Assume that $x_0,r,$ and $\lambda$ are supplied to the algorithm.
For $q\geq1$, $\eps\in(0,1/100)$, and an internal order
$Q=\max\{q,\widetilde O(1)\}$ meeting
\eqref{eq:cap-conditions}, there is a restart-until-success algorithm
returning $Y$ with
$\mathcal R_q(\law(Y)\|\pi)\leq\eps$ after
\begin{equation}\label{eq:cold-main}
  \widetilde O\!\left(
    \frac{\lambda}{r^2}
      \bigl(d^{5/2}\sqrt Q+Qd^2\bigr)
  \right)
\end{equation}
expected membership queries.  This statement uses only Klartag's
established $\mathcal C_d=O(\log(ed))$ bound.  Letwin's theorem merely
removes the hidden factor $\mathcal C_d^{3/2}$ from the annealing term.
\end{theorem}

\begin{proof}
The Poincar\'e part of Corollary~\ref{cor:scalar-heat} supplies the
known bound
\[
  \overline H=196\mathcal C_d^3d\lambda^2
  =O\!\left(d\lambda^2\log^3(ed)\right),
  \qquad \overline\Lambda=\lambda.
\]
Substitution into Theorem~\ref{thm:algorithm-main} gives the displayed
two terms, up to the factor $\mathcal C_d^{3/2}$ absorbed by
$\widetilde O$, as well as an additive $d$.

To absorb the latter, translate $x_0$ to zero and let
$\rho_K(\theta)=\sup\{t\geq0:t\theta\in K\}$ be the radial function of
$K$.  Since $B(0,r)\subseteq K$,
$\rho_K(\theta)\geq r$ for every $\theta\in S^{d-1}$.  Polar integration
gives
\[
  d\lambda=\E_\pi\norm X^2
  =\frac{d}{d+2}
    \frac{\int_{S^{d-1}}\rho_K(\theta)^{d+2}\dd\theta}
         {\int_{S^{d-1}}\rho_K(\theta)^d\dd\theta}
  \geq\frac{d}{d+2}r^2.
\]
Hence $\lambda/r^2\geq1/(d+2)$, and, since $Q\geq2$, the term
$Qd^2\lambda/r^2$ dominates $d$ up to a universal constant.  This
proves \eqref{eq:cold-main} without invoking Theorem~\ref{thm:letwin}.
Using that theorem instead sets $\overline H=6160d\lambda^2$ and only
removes the indicated polylogarithmic loss.
\end{proof}

For $r=\lambda=1$ and $Q=\widetilde O(1)$,
Theorem~\ref{thm:cold-main} gives
$\widetilde O(d^{2.5})$ expected membership queries.  In the same
membership-oracle model and finite-order R\'enyi convention,
\citet[Corollary~1.6]{KV26zeroth} give
\[
  \widetilde O\!\left(
    qd^2R^{3/2}\Lambda^{1/4}
    +qd^2\Lambda
  \right).
\]
For a centered isotropic body, $R=\sqrt d$ and $\Lambda=1$, making the
warm-start generation term $\widetilde O(qd^{2.75})$.  On the
exact-centered scalar-covariance subclass considered here, our gain is
confined to this warm-start generation term; unlike their general
theorem, ours assumes that the annealing center is the exact mean and
the covariance is scalar.  The final
$\widetilde O(Qd^2)$ warm sampling phase is imported from their work.
Their Section~6 and Appendix~A already obtain a $d^{2.5}$ exponent for certain
structured families---including cubes, simplices, and bodies of
revolution---under operator-norm quadratic-tilt stability.  Theorem
\ref{thm:cold-main} instead applies to every centered isotropic convex
body satisfying the stated inscribed-ball normalization.

The operator-stability counterexample in
\citet[Proposition~6.3]{KV26zeroth} is compatible with our theorem.  It
shows that a single eigenvalue of $\Sigma_\beta$ can grow polynomially.
Our determinant argument controls
$\tr(\Sigma_\beta^2)+m_\beta^\top\Sigma_\beta m_\beta$ and hence the
variance of the scalar radial statistic; it does not claim a uniform
operator-norm bound.

\paragraph{Limitations.}
The center $x_0$ is assumed to be the exact mean of the target.  If it
is not, energy monotonicity only gives a bound relative to
$\E_\pi\norm{X-x_0}^2$, and the determinant initialization
$m_0=0$, $\Sigma_0=I$ is lost.  Likewise, scalar covariance can be
handled by a scalar rescaling, but arbitrary whitening changes a
spherical Gaussian tilt into an ellipsoidal one.  Establishing the
fully anisotropic bound
\[
  \sup_\beta\Var_{\pi_\beta}\norm{X-x_0}^2
  \stackrel{?}{=}O(R^2\Lambda)
\]
requires a weighted determinant or spectral potential not proved here.
Finally, \eqref{eq:cold-main} is an expected \emph{membership-query}
bound, not merely a count of accepted or proper transitions; the
rejection trials, caps, and full-attempt restarts are included.

\section{Evaluation-oracle logconcave sampling}
\label{sec:lift}

Let $V:\R^d\to\R\cup\{+\infty\}$ be lower semicontinuous and convex,
and let
\[
  \pi(\dd x)=Z^{-1}e^{-V(x)}\dd x.
\]
The structural statements in this subsection allow extended-real
potentials.  The algorithmic theorem below returns to finite-valued
$V$, matching the stated scope of the imported sampler propositions.
The exponential lift of \citet{KV25sampling,KV26zeroth} uses the convex
epigraph
\begin{equation}\label{eq:lift-body}
  \mathcal K=\{(x,t)\in\R^d\times\R:V(x)\leq dt\}.
\end{equation}
One evaluation of $V(x)$ implements one membership query to
$\mathcal K$.  For a point $x_0\in\R^d$ and $\beta\geq0$, define
\begin{equation}\label{eq:lift-anneal}
  P_\beta(\dd x\dd t)
  \propto
  e^{-\beta\norm{x-x_0}^2/2-dt}
  \1_{\mathcal K}(x,t)\dd x\dd t.
\end{equation}

\begin{proposition}[Exact marginal transfer]
\label{prop:lift-transfer}
The $x$-marginal of $P_\beta$ is
\[
  \pi_\beta(\dd x)
  \propto e^{-V(x)-\beta\norm{x-x_0}^2/2}\dd x.
\]
Conditionally on $X=x$,
\[
  T=\frac{V(x)+E}{d},
  \qquad E\sim\operatorname{Exp}(1),
\]
and this conditional kernel does not depend on $\beta$.  Therefore, for
all $a>1$ and $\beta,\beta'\geq0$ for which the divergences are finite,
\begin{equation}\label{eq:lift-renyi-equality}
  \mathcal R_a(P_\beta\|P_{\beta'})
  =\mathcal R_a(\pi_\beta\|\pi_{\beta'}).
\end{equation}
\end{proposition}

\begin{proof}
Integrating \eqref{eq:lift-anneal} over
$t\geq V(x)/d$ produces $d^{-1}e^{-V(x)}$ times the Gaussian tilt.
The displayed conditional law follows by shifting a rate-$d$
exponential random variable.  Both joint laws are obtained from their
$x$-marginals by applying the same Markov kernel $x\mapsto\law(T\mid x)$.
Their Radon--Nikodym derivative is therefore a function of $x$ alone,
and direct integration in \eqref{eq:renyi} proves
\eqref{eq:lift-renyi-equality}.
\end{proof}

\subsection{A fixed truncation with uniform mass}

Put $V_0=\min V$ and
\[
  L_{\pi,g}=\{x:V(x)-V_0\leq10d\}.
\]
Suppose $x_0\in L_{\pi,g}$, set
\[
  R^2=\E_\pi\norm{X-x_0}^2,
  \qquad l=\log(4e),
  \qquad B=\frac{V(x_0)}d+5+13l,
\]
and consider Kook and Vempala's fixed convex truncation
\begin{equation}\label{eq:fixed-lift-truncation}
  A=\mathcal K\cap
  \bigl\{(x,t):\norm{x-x_0}\leq Rl,\ t\leq B\bigr\}.
\end{equation}
After translating $x$ and the lift coordinate, this is the lifted
epigraph intersected with the cylinder
$B_{Rl}(0)\times[-21,13l-6]$ used by their algorithm
\citep[Equation~(5.1)]{KV26zeroth}.  The apparent lower cutoff is already
implied by the epigraph and $V(x_0)-V_0\leq10d$.

\begin{lemma}[Uniform mass of the fixed truncation]
\label{lem:uniform-truncation-mass}
For every $\beta\geq0$,
\begin{equation}\label{eq:uniform-truncation-mass}
  p_\beta:=P_\beta(A)\geq\frac12.
\end{equation}
\end{lemma}

\begin{proof}
Let
\[
  W_\beta(x)=V(x)+\frac\beta2\norm{x-x_0}^2,
  \qquad
  S=T+\frac{\beta}{2d}\norm{X-x_0}^2.
\]
The unit-Jacobian triangular map $(x,t)\mapsto(x,s)$ sends
\eqref{eq:lift-anneal} exactly to the standard exponential lift
\[
  e^{-ds}\1_{\{W_\beta(x)\leq ds\}}\dd x\dd s
\]
of $W_\beta$.  Its $x$-marginal is $\pi_\beta$.  By energy monotonicity,
which here follows directly from
\[
  \frac{\dd}{\dd\beta}\E_{\pi_\beta}\norm{X-x_0}^2
  =-\frac12\Var_{\pi_\beta}\norm{X-x_0}^2\leq0,
\]
\begin{equation}\label{eq:lift-energy-monotonicity}
  R_\beta^2:=\E_{\pi_\beta}\norm{X-x_0}^2\leq R^2.
\end{equation}
Moreover,
\[
  W_\beta(x_0)-\min W_\beta
  \leq V(x_0)-V_0\leq10d,
  \qquad W_\beta(x_0)=V(x_0).
\]
We now reproduce the two tail estimates behind Kook and Vempala's
convex truncation.  Their logconcave exponential-decay lemma
\citep[Lemma~A.7]{KV25sampling} gives
\begin{equation}\label{eq:uniform-truncation-radial-tail}
  \Pp_{\pi_\beta}\!\left(\norm{X-x_0}>R_\beta l\right)
  \leq e^{-l+1}=\frac14.
\end{equation}
For completeness, the mean of $S$ admits a short direct bound.  Put
$m_\beta=\min W_\beta$ and
$F(a)=\int_{\R^d}e^{-a(W_\beta(x)-m_\beta)}\dd x$.  If $x_\beta$
minimizes $W_\beta$, convexity gives, for $a\geq1$,
\[
  a\left(W_\beta\!\left(x_\beta+
    \frac{z-x_\beta}{a}\right)-m_\beta\right)
  \leq W_\beta(z)-m_\beta.
\]
With $x=x_\beta+(z-x_\beta)/a$ and
$\dd x=a^{-d}\dd z$, this yields
$F(a)\geq a^{-d}F(1)$.  For $h>0$, it follows that
\[
  \frac{F(1)-F(1+h)}h
  \leq\frac{1-(1+h)^{-d}}hF(1).
\]
The integrand on the left is
$e^{-U}(1-e^{-hU})/h$ for $U=W_\beta-m_\beta\geq0$; monotone
convergence as $h\downarrow0$ therefore gives
$\E_{\pi_\beta}(W_\beta-m_\beta)\leq d$.  Since conditionally on $X$,
$S=(W_\beta(X)+E)/d$ for $E\sim\operatorname{Exp}(1)$,
\begin{equation}\label{eq:uniform-truncation-lift-mean}
  \E S\leq\frac{m_\beta}{d}+1+\frac1d
  \leq\frac{W_\beta(x_0)}d+2.
\end{equation}
Kook and Vempala prove $\Var(S)\leq160$ for every such exponential
lift and record the needed exponential-decay lemma
\citep[Lemmas~2.5 and~A.7]{KV25sampling}.  The $S$-marginal is
logconcave, so applying that lemma after centering and using
$\sqrt{160}<13$ gives
\begin{equation}\label{eq:uniform-truncation-lift-tail}
  \Pp\!\left(S>\frac{W_\beta(x_0)}d+5+13l\right)
  \leq\Pp(S-\E S>13l)
  \leq e^{-l+1}=\frac14.
\end{equation}
Thus a union bound assigns mass at least $1/2$ to
\[
  \left\{\norm{x-x_0}\leq R_\beta l,\quad
  s\leq\frac{W_\beta(x_0)}d+5+13l\right\}.
\]
Under the inverse triangular map this event lies in $A$: the radial
bound follows from \eqref{eq:lift-energy-monotonicity}, the upper cutoff
is $B$, and
$t=s-\beta\norm{x-x_0}^2/(2d)\leq s$.  This proves
\eqref{eq:uniform-truncation-mass}.
\end{proof}

Let $\overline P_\beta=P_\beta(\cdot\mid A)$ and let $\nu_\beta$ be its
$x$-marginal.  Because $A$ is fixed,
$\nu_\beta(\dd x)\propto e^{-\beta\norm{x-x_0}^2/2}\nu_0(\dd x)$.
The preceding lemma transfers the exact heat bound to the implemented
path without requiring the truncated marginal itself to be isotropic.

\begin{proposition}[Heat and overlap on the truncated path]
\label{prop:truncated-path-transfer}
For $U(x)=\norm{x-x_0}^2$ and every $\beta\geq0$,
\begin{equation}\label{eq:truncated-heat-transfer}
  \Var_{\nu_\beta}U
  \leq2\Var_{\pi_\beta}U.
\end{equation}
Moreover, for $a>1$ and $\beta,\beta'\geq0$,
\begin{equation}\label{eq:truncated-renyi-equality}
  \mathcal R_a(\overline P_\beta\|\overline P_{\beta'})
  =\mathcal R_a(\nu_\beta\|\nu_{\beta'}),
\end{equation}
and
\begin{equation}\label{eq:zero-truncation-warm}
  \mathcal R_\infty(\overline P_0\|P_0)
  =\log\frac1{p_0}\leq\log2.
\end{equation}
\end{proposition}

\begin{proof}
Conditioning and marginalization give
$\dd\nu_\beta/\dd\pi_\beta\leq p_\beta^{-1}\leq2$.  Hence
\[
  \Var_{\nu_\beta}U
  \leq\E_{\nu_\beta}(U-\E_{\pi_\beta}U)^2
  \leq2\E_{\pi_\beta}(U-\E_{\pi_\beta}U)^2,
\]
which is \eqref{eq:truncated-heat-transfer}.  On the common fixed
support $A$, the conditional density of $T$ given $X=x$ is proportional
to
\[
  e^{-dt}\1_{\{V(x)/d\leq t\leq B\}},
\]
independently of $\beta$.  The proof of
Proposition~\ref{prop:lift-transfer} therefore gives
\eqref{eq:truncated-renyi-equality}.  Finally,
$\overline P_0=P_0(\cdot\mid A)$, so
\eqref{eq:zero-truncation-warm} is immediate from
Lemma~\ref{lem:uniform-truncation-mass}.
\end{proof}

\subsection{The cold-start algorithm}

We now prove the formal version of
Theorem~\ref{thm:intro-evaluation}.  Its assumptions match the
well-rounded evaluation-oracle regime of Kook and Vempala.

\begin{theorem}[Evaluation-oracle cold start]
\label{thm:evaluation-cold-main}
Let $V:\R^d\to\R$ be convex and given by an evaluation oracle.  Suppose
$\pi(\dd x)\propto e^{-V(x)}\dd x$ is centered and isotropic and
$B(0,1)\subseteq L_{\pi,g}$.  Fix $q\geq1$ and
$\eps\in(0,1/100)$.  For a sufficiently large universal constant
$C_0$, take
\begin{equation}\label{eq:evaluation-order-condition}
  Q=\max\left\{q,2,
    \left\lceil C_0\log^2\!\left(e^2+d+q+\frac1\eps\right)\right\rceil
    \right\}.
\end{equation}
There is a restart-until-success algorithm whose output $Y$ satisfies
\[
  \mathcal R_q(\law(Y)\|\pi)\leq\eps
\]
and whose expected evaluation-query complexity is
\begin{equation}\label{eq:evaluation-query-bound}
  \widetilde O\!\left(d^{5/2}\sqrt Q+Qd^2\right).
\end{equation}
Here $\widetilde O$ suppresses logarithms in $d,Q,$ and $1/\eps$; no
additional polynomial dependence on $Q$ is hidden.
No use of Theorem~\ref{thm:letwin} is required.
\end{theorem}

\begin{proof}
We first describe an ideal uncapped attempt and verify its R\'enyi
invariants.  Rejection caps and full-attempt restarts are added at the
end.  Every hatted law below refers to this ideal attempt; the query
costs and cap-hit probabilities refer to the capped implementations
coupled to the same ideal intermediate laws.

We use the fixed set $A$ in \eqref{eq:fixed-lift-truncation} with
$x_0=0$.  Isotropy gives $R=\sqrt d$, so its radial scale
$D=Rl$ is $\Theta(\sqrt d)$.  In the notation of the imported
construction, its annealing targets are
\begin{equation}\label{eq:KV-lift-targets}
  \mu_{\sigma^2,\rho}(\dd x\dd t)
  \propto
  e^{-\norm x^2/(2\sigma^2)-\rho t}\1_A(x,t)\dd x\dd t.
\end{equation}
When $\rho=d$ and $\beta=1/\sigma^2$, this is exactly
$\overline P_\beta$.

Invoke the initialization and first two
ideal phases of the R\'enyi warm-start construction of Kook and Vempala
\citep[Section~5.2]{KV26zeroth}.  Choose their preserved base order
\[
  q_0=\max\left\{2,
    \left\lceil C_1\log\!\left(e+d+Q+\frac1\eps\right)\right\rceil
  \right\},
\]
where $C_1$ is a sufficiently large universal constant.  Increasing
$C_0$ if necessary ensures $q_0\leq Q$ and all rejection-cap
conditions below.  Their Phase~I costs
$\widetilde O(q_0^{1/2}d^{2.5})$, and their Phase~II costs
\[
  \widetilde O\!\left(q_0^{1/2}d^{2.5}+q_0d^2D\right)
  =\widetilde O(d^{2.5}).
\]
Cap the final inner update at $\sigma^2=1$; this only decreases its
parameter change, so the same neighboring-target overlap estimate
applies.
Thus these phases produce a law $\widehat P_1$ satisfying
\begin{equation}\label{eq:phase-two-invariant}
  \mathcal R_{2q_0}(\widehat P_1\|\overline P_1)\leq c_0
\end{equation}
with $c_0=1/10$.  At this point the
lift coefficient is $\rho=d$, and the exact target is precisely
$\overline P_1$.

By Theorem~\ref{thm:poincare-heat}, Klartag's
$\mathcal C_d=O(\log(ed))$ estimate, and
Proposition~\ref{prop:truncated-path-transfer}, one may fix a universal
certified constant $C_H$ such that the truncated path has the numerical
heat bound
\begin{equation}\label{eq:truncated-Hbar}
  \overline H:=C_Hd\log^3(ed),
  \qquad
  \Var_{\nu_\beta}\norm X^2\leq\overline H
  \quad(\beta\geq0).
\end{equation}
Set
\[
  \Delta=\min\left\{1,\frac1{4\sqrt{Q\overline H}}\right\},
  \qquad L=\left\lceil\frac1\Delta\right\rceil,
  \qquad h=\frac1L,
\]
and use the positive precisions $1,(L-1)h,\ldots,h$.  Applying
Lemma~\ref{lem:partition-identities} to the base marginal of the
truncated path, and then using
\eqref{eq:truncated-renyi-equality}, shows that consecutive exact joint
targets satisfy
\begin{equation}\label{eq:truncated-neighbor-overlap}
  \mathcal R_{2q_0}(\overline P_{kh}\|\overline P_{(k-1)h})
  \leq\frac{2q_0}{8}h^2\overline H
  \leq\frac1{64}
\end{equation}
for $k=2,\ldots,L$.

At each new positive precision, run Kook and Vempala's truncated
proximal sampler $\mathsf{PS}_{\rm ann}$ from order $q_0$ to order
$2q_0$ and restore the invariant
$\mathcal R_{2q_0}\leq c_0$.  The weak triangle inequality and
\eqref{eq:truncated-neighbor-overlap} keep its input warmness below
$10$.  Proposition~5.2 of \citet{KV26zeroth} charges, at target
$\beta=1/\sigma^2$,
\begin{equation}\label{eq:truncated-phase-cost}
  \widetilde O\!\left(d^2(\sigma^2\vee1)\right)
  =\widetilde O\!\left(d^2(\beta^{-1}\vee1)\right)
\end{equation}
evaluation queries.  Consequently all positive-precision calls cost
\begin{equation}\label{eq:truncated-harmonic-cost}
  \sum_{k=1}^{L-1}\widetilde O\!\left(\frac{d^2}{kh}\right)
  =\widetilde O\!\left(\frac{d^2}{h}\right)
  =\widetilde O\!\left(d^2\sqrt{Q\overline H}\right).
\end{equation}

At the last positive target $\overline P_h$, make one additional
same-target call to $\mathsf{PS}_{\rm ann}$, now from order $q_0$ to
order $2Q$, obtaining a law $\widehat P_h$ with
$\mathcal R_{2Q}(\widehat P_h\|\overline P_h)\leq c_0$.  The target-order
dependence of this call is only logarithmic, and its cost is absorbed by
\eqref{eq:truncated-harmonic-cost}.  The terminal exact overlap is
\[
  \mathcal R_{2Q}(\overline P_h\|\overline P_0)
  \leq\frac{2Q}{8}h^2\overline H\leq\frac1{64}.
\]
The weak triangle inequality therefore gives
\begin{equation}\label{eq:full-lift-warm-input}
  \mathcal R_Q(\widehat P_h\|\overline P_0)<1
\end{equation}
because $2c_0+1/64<1$.

Let $\rho$ be the $x$-marginal of $\widehat P_h$ and let $\nu_0$ be the
$x$-marginal of $\overline P_0$.  By data processing,
$M_Q(\rho\|\nu_0)<e$.  Also
$\dd\nu_0/\dd\pi\leq p_0^{-1}\leq2$, and hence
\begin{equation}\label{eq:exact-target-warmness}
  M_Q(\rho\|\pi)
  \leq2^{(Q-1)/Q}M_Q(\rho\|\nu_0)<2e<10.
\end{equation}
Resample $T=(V(X)+E)/d$ conditionally on $X\sim\rho$.  Because this is
the same conditional kernel as under $P_0$, the resulting joint
warmness relative to $P_0$ is exactly $M_Q(\rho\|\pi)<10$.  Now run the
full, untruncated exponential-lift sampler $\mathsf{PS}_{\rm exp}$ of
Kook and Vempala.  Its balanced finite-order guarantee
\citep[Proposition~5.1]{KV26zeroth}, together with isotropy
$\norm{\Cov(\pi)}_{\rm op}=1$, produces order-$Q$ error at most
$\eps/2$ using $\widetilde O(Qd^2)$ additional evaluation queries.

Put back the capped implementations coupled to the ideal attempt just
analyzed.  Give the calls total cap budget $\eta=\eps/8$, and restart the
entire pipeline from a fresh initialization if any cap is hit.  The phase
count and all iteration horizons have logarithms bounded
by a universal multiple of
$\log(e^2+d+Q+1/\eps)$.  Thus one may choose $q_0$ polylogarithmically
and then take $C_0$ in \eqref{eq:evaluation-order-condition} large
enough to meet every cap condition.  The unconditional expected cost of
one capped attempt is bounded by the sum above, because a call is made
only if all earlier calls succeeded.  A union bound makes one attempt
successful with probability at least $1-\eta$.  Let $\rho_*$ be the
ideal final $x$-law, let $S$ be the event that no cap is hit, and let
$\nu=\law(Y\mid S)$ be the law returned by a successful attempt.  The
coupling gives
\[
  \frac{\dd\nu}{\dd\rho_*}\leq\frac1{\Pp(S)}\leq\frac1{1-\eta}.
\]
Directly from the definition of R\'enyi divergence,
\[
  \mathcal R_Q(\nu\|\pi)
  \leq\mathcal R_Q(\rho_*\|\pi)
     +\frac{Q}{Q-1}\log\frac1{1-\eta}
  \leq\eps.
\]
Here the ideal final sampler was run to error at most $\eps/2$,
$Q\geq2$, $\eta=\eps/8$, and $\eps<1/100$.  Full restarts increase the
expected cost by only a constant factor and terminate almost surely.

Finally, \eqref{eq:truncated-Hbar} turns
\eqref{eq:truncated-harmonic-cost} into
$\widetilde O(d^{5/2}\sqrt Q)$.  Adding Phases~I--II and the final
balanced call proves \eqref{eq:evaluation-query-bound}; monotonicity in
the R\'enyi order gives the requested order $q$.
\end{proof}

For comparison, the evaluation-oracle bound of
\citet[Corollary~1.9]{KV26zeroth} is
\[
  \widetilde O\!\left(
    d^{2.5}+qd^2R^{3/2}(\Lambda^{1/4}\vee1)
    +qd^2(\Lambda\vee1)
  \right),
\]
up to logarithmic accuracy factors, where
$R^2=\E_\pi\norm X^2$ and
$\Lambda=\norm{\Cov(\pi)}_{\rm op}$.  In centered isotropic position
this is $\widetilde O(qd^{2.75})$.  Theorem
\ref{thm:evaluation-cold-main} replaces the third-phase term on exactly
that subclass; it retains their ground-set normalization, preprocessing,
truncated sampler, and final full-lift sampler.  The later unified
analysis of \citet{KV26unified} improves the last warm-start sampling
term, but does not supply the order-boosting truncated sampler used in
the annealing phases.  As throughout the paper, exact centering and
isotropy are assumptions of the result, not outputs of a rounding
procedure.

\begingroup
\raggedright
\bibliographystyle{plainnat}
\bibliography{references}
\endgroup

\end{document}
